%% annex/p08/temporal-dynamic-epistemic.tex --- Part 8.3 proof sketches
%% Both items inherit sketch status from v1.1-verified; neither is ratio-tuned.

\begin{annexitem}{annex:p8.temporal.1}{thm:p08-ltl-adequacy}{LTL soundness and weak completeness}
\sketchnote{the consistency proof for the maximally consistent closure sets;
the eventuality-realization lemma for every pending strong-until obligation;
and the full truth-lemma case calculations.}
For soundness, propositional instances and modus ponens preserve validity.
The two next-time principles use the unique successor of each position.
Strong-until unfolding separates a witness at the current position from a
strictly later witness.  For the induction rule, suppose its theorem-only
premise is valid and let \(m\) be the least witness for the right operand of an
until formula.  The premise gives the target at \(m\), and backward induction
through the intervening positions gives it at the starting position.  Thus
every rule on \hyperref[sheet:hilbert-ltl]{\texttt{sheet:hilbert-ltl}}
preserves infinite-word validity.

For completeness, close the underivable formula under subformulas, negation,
and the obligations generated by strong until.  The omitted consistency
argument supplies maximally consistent sets over this finite closure.  Link a
set to a successor set according to its \(\mathsf X\)-requirements.  The omitted
eventuality-realization lemma shows that every pending until obligation obtains
a witness along a suitable linked path.  Unravelling that path gives an
infinite word.  The omitted structural truth lemma identifies membership in
each closure set with truth at the corresponding position, so an underivable
formula has a falsifying word and position.  This proves only the weak
theoremhood statement of Theorem~\ref{thm:p08-ltl-adequacy}; arbitrary-premise
strong completeness and finite-trace semantics are outside the claim.
\end{annexitem}

\begin{annexitem}{annex:p8.temporal.2}{thm:p08-s5n-adequacy}{Multi-agent \(\mathsf{S5}_{I}\) adequacy}
\sketchnote{the Lindenbaum extension from a nonderivation to a maximally
consistent set; the modal existence lemma for each agent-indexed canonical
relation; and the full structural truth-lemma calculation.}
Soundness is componentwise.  Distribution and theorem-only necessitation are
valid for every relation; \(T_i\), \(4_i\), and \(5_i\) use respectively
reflexivity, transitivity, and Euclideanness.  Equivalence relations have all
three properties, and no soundness step relates two distinct agents.

For completeness, take maximally consistent sets as worlds and define
\(w\sim_i v\) exactly when every \(K_iA\in w\) has \(A\in v\).  The three
non-distribution schemas make this canonical relation reflexive, transitive,
and Euclidean, hence an equivalence relation.  If
\(\Gamma\not\derives_{\mathsf{S5}_{I}}A\), the omitted Lindenbaum argument
produces a world containing \(\Gamma\) and \(\neg A\).  The omitted modal
existence lemma supplies the successor required in the negative \(K_i\)-case,
and the omitted truth lemma then makes every member of \(\Gamma\) true there
while \(A\) is false.  This yields the local-consequence direction in
Theorem~\ref{thm:p08-s5n-adequacy}, again without any cross-agent interaction
principle.
\end{annexitem}
